% =====================================================================
% Note on: Bond Total Return Swaps — Theory, Pricing & Practice
% Nicholas Burgess — ssrn author 1728976
% =====================================================================
\documentclass[11pt,a4paper]{article}

\newcommand{\papertag}{burgess\_trs}

% =====================================================================
% Shared preamble for paper notes
% Include with:  % =====================================================================
% Shared preamble for paper notes
% Include with:  % =====================================================================
% Shared preamble for paper notes
% Include with:  \input{../_common/preamble}
% =====================================================================

% --- Core packages ---
\usepackage[utf8]{inputenc}
\usepackage[T1]{fontenc}
\usepackage{lmodern}
\usepackage[margin=2.5cm]{geometry}
\usepackage{amsmath,amssymb,amsthm,mathtools}
\usepackage{bm}
\usepackage{booktabs}
\usepackage{array}
\usepackage{enumitem}
\usepackage{hyperref}
\usepackage{xcolor}
\usepackage{listings}
\usepackage{fancyhdr}
\usepackage{titlesec}

% --- Hyperref styling ---
\hypersetup{
  colorlinks=true,
  linkcolor=blue!50!black,
  urlcolor=blue!50!black,
  citecolor=blue!50!black
}

% --- Theorem-like environments ---
\theoremstyle{definition}
\newtheorem{defn}{Definition}
\newtheorem{remark}{Remark}
\newtheorem{example}{Example}

\theoremstyle{plain}
\newtheorem{prop}{Proposition}
\newtheorem{lemma}{Lemma}

% --- Coloured boxes for the structured sections ---
% Validation block, commentary, TODOs use coloured frames so the eye
% can find them when flipping through the note.
\usepackage[most]{tcolorbox}

\newtcolorbox{commentary}{
  colback=yellow!5,
  colframe=yellow!50!black,
  title=Commentary,
  fonttitle=\bfseries,
  breakable
}

\newtcolorbox{validation}{
  colback=green!3,
  colframe=green!50!black,
  title=Validation,
  fonttitle=\bfseries,
  breakable
}

\newtcolorbox{todobox}{
  colback=red!3,
  colframe=red!50!black,
  title=TODO / Open questions,
  fonttitle=\bfseries,
  breakable
}

% --- Code listing style for Python snippets in validation blocks ---
\lstdefinestyle{py}{
  language=Python,
  basicstyle=\ttfamily\footnotesize,
  keywordstyle=\color{blue!60!black},
  commentstyle=\color{gray},
  stringstyle=\color{red!60!black},
  showstringspaces=false,
  breaklines=true,
  frame=single,
  framesep=4pt,
  rulecolor=\color{gray!30}
}

% --- Page header: shows paper tag so a printed note is identifiable ---
\pagestyle{fancy}
\fancyhf{}
\fancyhead[L]{\small\textit{\papertag}}
\fancyhead[R]{\small\thepage}
\renewcommand{\headrulewidth}{0.4pt}

% --- Section spacing: tighter than article default ---
\titlespacing*{\section}{0pt}{1.5ex plus 0.5ex minus 0.2ex}{1ex plus 0.2ex}
\titlespacing*{\subsection}{0pt}{1.2ex plus 0.3ex minus 0.2ex}{0.6ex plus 0.2ex}

% =====================================================================

% --- Core packages ---
\usepackage[utf8]{inputenc}
\usepackage[T1]{fontenc}
\usepackage{lmodern}
\usepackage[margin=2.5cm]{geometry}
\usepackage{amsmath,amssymb,amsthm,mathtools}
\usepackage{bm}
\usepackage{booktabs}
\usepackage{array}
\usepackage{enumitem}
\usepackage{hyperref}
\usepackage{xcolor}
\usepackage{listings}
\usepackage{fancyhdr}
\usepackage{titlesec}

% --- Hyperref styling ---
\hypersetup{
  colorlinks=true,
  linkcolor=blue!50!black,
  urlcolor=blue!50!black,
  citecolor=blue!50!black
}

% --- Theorem-like environments ---
\theoremstyle{definition}
\newtheorem{defn}{Definition}
\newtheorem{remark}{Remark}
\newtheorem{example}{Example}

\theoremstyle{plain}
\newtheorem{prop}{Proposition}
\newtheorem{lemma}{Lemma}

% --- Coloured boxes for the structured sections ---
% Validation block, commentary, TODOs use coloured frames so the eye
% can find them when flipping through the note.
\usepackage[most]{tcolorbox}

\newtcolorbox{commentary}{
  colback=yellow!5,
  colframe=yellow!50!black,
  title=Commentary,
  fonttitle=\bfseries,
  breakable
}

\newtcolorbox{validation}{
  colback=green!3,
  colframe=green!50!black,
  title=Validation,
  fonttitle=\bfseries,
  breakable
}

\newtcolorbox{todobox}{
  colback=red!3,
  colframe=red!50!black,
  title=TODO / Open questions,
  fonttitle=\bfseries,
  breakable
}

% --- Code listing style for Python snippets in validation blocks ---
\lstdefinestyle{py}{
  language=Python,
  basicstyle=\ttfamily\footnotesize,
  keywordstyle=\color{blue!60!black},
  commentstyle=\color{gray},
  stringstyle=\color{red!60!black},
  showstringspaces=false,
  breaklines=true,
  frame=single,
  framesep=4pt,
  rulecolor=\color{gray!30}
}

% --- Page header: shows paper tag so a printed note is identifiable ---
\pagestyle{fancy}
\fancyhf{}
\fancyhead[L]{\small\textit{\papertag}}
\fancyhead[R]{\small\thepage}
\renewcommand{\headrulewidth}{0.4pt}

% --- Section spacing: tighter than article default ---
\titlespacing*{\section}{0pt}{1.5ex plus 0.5ex minus 0.2ex}{1ex plus 0.2ex}
\titlespacing*{\subsection}{0pt}{1.2ex plus 0.3ex minus 0.2ex}{0.6ex plus 0.2ex}

% =====================================================================

% --- Core packages ---
\usepackage[utf8]{inputenc}
\usepackage[T1]{fontenc}
\usepackage{lmodern}
\usepackage[margin=2.5cm]{geometry}
\usepackage{amsmath,amssymb,amsthm,mathtools}
\usepackage{bm}
\usepackage{booktabs}
\usepackage{array}
\usepackage{enumitem}
\usepackage{hyperref}
\usepackage{xcolor}
\usepackage{listings}
\usepackage{fancyhdr}
\usepackage{titlesec}

% --- Hyperref styling ---
\hypersetup{
  colorlinks=true,
  linkcolor=blue!50!black,
  urlcolor=blue!50!black,
  citecolor=blue!50!black
}

% --- Theorem-like environments ---
\theoremstyle{definition}
\newtheorem{defn}{Definition}
\newtheorem{remark}{Remark}
\newtheorem{example}{Example}

\theoremstyle{plain}
\newtheorem{prop}{Proposition}
\newtheorem{lemma}{Lemma}

% --- Coloured boxes for the structured sections ---
% Validation block, commentary, TODOs use coloured frames so the eye
% can find them when flipping through the note.
\usepackage[most]{tcolorbox}

\newtcolorbox{commentary}{
  colback=yellow!5,
  colframe=yellow!50!black,
  title=Commentary,
  fonttitle=\bfseries,
  breakable
}

\newtcolorbox{validation}{
  colback=green!3,
  colframe=green!50!black,
  title=Validation,
  fonttitle=\bfseries,
  breakable
}

\newtcolorbox{todobox}{
  colback=red!3,
  colframe=red!50!black,
  title=TODO / Open questions,
  fonttitle=\bfseries,
  breakable
}

% --- Code listing style for Python snippets in validation blocks ---
\lstdefinestyle{py}{
  language=Python,
  basicstyle=\ttfamily\footnotesize,
  keywordstyle=\color{blue!60!black},
  commentstyle=\color{gray},
  stringstyle=\color{red!60!black},
  showstringspaces=false,
  breaklines=true,
  frame=single,
  framesep=4pt,
  rulecolor=\color{gray!30}
}

% --- Page header: shows paper tag so a printed note is identifiable ---
\pagestyle{fancy}
\fancyhf{}
\fancyhead[L]{\small\textit{\papertag}}
\fancyhead[R]{\small\thepage}
\renewcommand{\headrulewidth}{0.4pt}

% --- Section spacing: tighter than article default ---
\titlespacing*{\section}{0pt}{1.5ex plus 0.5ex minus 0.2ex}{1ex plus 0.2ex}
\titlespacing*{\subsection}{0pt}{1.2ex plus 0.3ex minus 0.2ex}{0.6ex plus 0.2ex}

% =====================================================================
% House notation: shared symbols used across all paper notes.
% Each note imports this and may shadow individual macros if a paper
% uses a clashing convention — but the *default* meaning is fixed here.
%
% Include with:  % =====================================================================
% House notation: shared symbols used across all paper notes.
% Each note imports this and may shadow individual macros if a paper
% uses a clashing convention — but the *default* meaning is fixed here.
%
% Include with:  % =====================================================================
% House notation: shared symbols used across all paper notes.
% Each note imports this and may shadow individual macros if a paper
% uses a clashing convention — but the *default* meaning is fixed here.
%
% Include with:  \input{../_common/notation}
% =====================================================================

% --- Measures and expectations ---
\newcommand{\Q}{\mathbb{Q}}              % risk-neutral measure
\newcommand{\Qt}{\mathbb{Q}^{T}}         % T-forward measure
\newcommand{\E}{\mathbb{E}}              % expectation
\newcommand{\Et}[1]{\E^{#1}}             % expectation under measure #1
\newcommand{\Ftime}[1]{\mathcal{F}_{#1}} % filtration at #1

% --- Discount and forward objects ---
\newcommand{\Disc}[2]{D_{#1,#2}}         % discount factor D_{t,T}
\newcommand{\Bond}[2]{P_{#1}(#2)}        % bond price P_t(y)
\newcommand{\Ann}[1]{A_{#1}}             % annuity / level
\newcommand{\Numer}[1]{N_{#1}}           % numeraire
\newcommand{\Fwd}[2]{F_{#1,#2}}          % forward price F_{t,T}
\newcommand{\Fwdpx}[2]{\mathrm{Fwd}_{#1}(#2)} % verbose forward-price F_t(T)

% --- Rates ---
\newcommand{\IRR}{\mathrm{IRR}}          % internal rate of return
\newcommand{\Yld}{y}                     % bond yield variable
\newcommand{\Strike}{K}                  % strike (price-space)
\newcommand{\Lock}{L}                    % locked yield / rate
\newcommand{\Repo}{r^{\mathrm{repo}}}    % repo rate
\newcommand{\Fund}{r^{\mathrm{fun}}}     % unsecured funding rate
\newcommand{\OIS}{r^{\mathrm{ois}}}      % OIS rate
\newcommand{\Coll}{r^{\mathrm{c}}}       % collateral rate
\newcommand{\Rcmt}{R^{\mathrm{cmt}}}     % constant-maturity-treasury rate
\newcommand{\Rswp}{R^{\mathrm{swp}}}     % swap rate (used for risky variant)
\newcommand{\Rec}{\mathrm{Rec}}          % recovery rate
\newcommand{\NPV}{\mathrm{NPV}}          % present-value functional
\newcommand{\PVone}{\mathrm{PV01}}       % risky annuity
\newcommand{\Loss}{\mathrm{Loss}}        % default-leg integral
\newcommand{\Sasw}{S^{\mathrm{ASW}}}     % asset-swap spread
\newcommand{\Scds}{S^{\mathrm{CDS}}}     % CDS spread
\newcommand{\Basis}{\mathrm{Basis}}      % CDS-bond basis

% --- Sensitivities ---
\newcommand{\Diff}[2]{D_{#1}^{#2}}       % Euler's notation: D_x^k[f]
\newcommand{\Riskf}{\mathrm{RiskFactor}} % bond risk factor (= -DV01*1e4)
\newcommand{\DVone}{\mathrm{DV01}}

% --- Indicator / misc ---
\newcommand{\Ind}{\mathbf{1}}
\newcommand{\given}{\,|\,}
\newcommand{\dd}{\,\mathrm{d}}

% --- Greeks-style ---
\newcommand{\Gam}{\Gamma}
\newcommand{\Vega}{\mathcal{V}}

% =====================================================================

% --- Measures and expectations ---
\newcommand{\Q}{\mathbb{Q}}              % risk-neutral measure
\newcommand{\Qt}{\mathbb{Q}^{T}}         % T-forward measure
\newcommand{\E}{\mathbb{E}}              % expectation
\newcommand{\Et}[1]{\E^{#1}}             % expectation under measure #1
\newcommand{\Ftime}[1]{\mathcal{F}_{#1}} % filtration at #1

% --- Discount and forward objects ---
\newcommand{\Disc}[2]{D_{#1,#2}}         % discount factor D_{t,T}
\newcommand{\Bond}[2]{P_{#1}(#2)}        % bond price P_t(y)
\newcommand{\Ann}[1]{A_{#1}}             % annuity / level
\newcommand{\Numer}[1]{N_{#1}}           % numeraire
\newcommand{\Fwd}[2]{F_{#1,#2}}          % forward price F_{t,T}
\newcommand{\Fwdpx}[2]{\mathrm{Fwd}_{#1}(#2)} % verbose forward-price F_t(T)

% --- Rates ---
\newcommand{\IRR}{\mathrm{IRR}}          % internal rate of return
\newcommand{\Yld}{y}                     % bond yield variable
\newcommand{\Strike}{K}                  % strike (price-space)
\newcommand{\Lock}{L}                    % locked yield / rate
\newcommand{\Repo}{r^{\mathrm{repo}}}    % repo rate
\newcommand{\Fund}{r^{\mathrm{fun}}}     % unsecured funding rate
\newcommand{\OIS}{r^{\mathrm{ois}}}      % OIS rate
\newcommand{\Coll}{r^{\mathrm{c}}}       % collateral rate
\newcommand{\Rcmt}{R^{\mathrm{cmt}}}     % constant-maturity-treasury rate
\newcommand{\Rswp}{R^{\mathrm{swp}}}     % swap rate (used for risky variant)
\newcommand{\Rec}{\mathrm{Rec}}          % recovery rate
\newcommand{\NPV}{\mathrm{NPV}}          % present-value functional
\newcommand{\PVone}{\mathrm{PV01}}       % risky annuity
\newcommand{\Loss}{\mathrm{Loss}}        % default-leg integral
\newcommand{\Sasw}{S^{\mathrm{ASW}}}     % asset-swap spread
\newcommand{\Scds}{S^{\mathrm{CDS}}}     % CDS spread
\newcommand{\Basis}{\mathrm{Basis}}      % CDS-bond basis

% --- Sensitivities ---
\newcommand{\Diff}[2]{D_{#1}^{#2}}       % Euler's notation: D_x^k[f]
\newcommand{\Riskf}{\mathrm{RiskFactor}} % bond risk factor (= -DV01*1e4)
\newcommand{\DVone}{\mathrm{DV01}}

% --- Indicator / misc ---
\newcommand{\Ind}{\mathbf{1}}
\newcommand{\given}{\,|\,}
\newcommand{\dd}{\,\mathrm{d}}

% --- Greeks-style ---
\newcommand{\Gam}{\Gamma}
\newcommand{\Vega}{\mathcal{V}}

% =====================================================================

% --- Measures and expectations ---
\newcommand{\Q}{\mathbb{Q}}              % risk-neutral measure
\newcommand{\Qt}{\mathbb{Q}^{T}}         % T-forward measure
\newcommand{\E}{\mathbb{E}}              % expectation
\newcommand{\Et}[1]{\E^{#1}}             % expectation under measure #1
\newcommand{\Ftime}[1]{\mathcal{F}_{#1}} % filtration at #1

% --- Discount and forward objects ---
\newcommand{\Disc}[2]{D_{#1,#2}}         % discount factor D_{t,T}
\newcommand{\Bond}[2]{P_{#1}(#2)}        % bond price P_t(y)
\newcommand{\Ann}[1]{A_{#1}}             % annuity / level
\newcommand{\Numer}[1]{N_{#1}}           % numeraire
\newcommand{\Fwd}[2]{F_{#1,#2}}          % forward price F_{t,T}
\newcommand{\Fwdpx}[2]{\mathrm{Fwd}_{#1}(#2)} % verbose forward-price F_t(T)

% --- Rates ---
\newcommand{\IRR}{\mathrm{IRR}}          % internal rate of return
\newcommand{\Yld}{y}                     % bond yield variable
\newcommand{\Strike}{K}                  % strike (price-space)
\newcommand{\Lock}{L}                    % locked yield / rate
\newcommand{\Repo}{r^{\mathrm{repo}}}    % repo rate
\newcommand{\Fund}{r^{\mathrm{fun}}}     % unsecured funding rate
\newcommand{\OIS}{r^{\mathrm{ois}}}      % OIS rate
\newcommand{\Coll}{r^{\mathrm{c}}}       % collateral rate
\newcommand{\Rcmt}{R^{\mathrm{cmt}}}     % constant-maturity-treasury rate
\newcommand{\Rswp}{R^{\mathrm{swp}}}     % swap rate (used for risky variant)
\newcommand{\Rec}{\mathrm{Rec}}          % recovery rate
\newcommand{\NPV}{\mathrm{NPV}}          % present-value functional
\newcommand{\PVone}{\mathrm{PV01}}       % risky annuity
\newcommand{\Loss}{\mathrm{Loss}}        % default-leg integral
\newcommand{\Sasw}{S^{\mathrm{ASW}}}     % asset-swap spread
\newcommand{\Scds}{S^{\mathrm{CDS}}}     % CDS spread
\newcommand{\Basis}{\mathrm{Basis}}      % CDS-bond basis

% --- Sensitivities ---
\newcommand{\Diff}[2]{D_{#1}^{#2}}       % Euler's notation: D_x^k[f]
\newcommand{\Riskf}{\mathrm{RiskFactor}} % bond risk factor (= -DV01*1e4)
\newcommand{\DVone}{\mathrm{DV01}}

% --- Indicator / misc ---
\newcommand{\Ind}{\mathbf{1}}
\newcommand{\given}{\,|\,}
\newcommand{\dd}{\,\mathrm{d}}

% --- Greeks-style ---
\newcommand{\Gam}{\Gamma}
\newcommand{\Vega}{\mathcal{V}}


\title{Note on: \emph{Bond Total Return Swaps --- Theory, Pricing \& Practice}\\
       \large Nicholas Burgess}
\author{Reading notes}
\date{\today}

\begin{document}
\maketitle

% =====================================================================
\section*{Metadata}
\begin{tabular}{@{}p{3.5cm}p{11cm}@{}}
\toprule
Title           & Bond Total Return Swaps: Theory, Pricing \& Practice \\
Author          & Nicholas Burgess \\
Date            & Not stated on cover \\
Source          & \href{https://ssrn.com/author=1728976}{ssrn:author=1728976};
                  companion repo \url{https://github.com/nburgessx/SwapsBook} \\
Local file      & \texttt{papers/burgess\_trs/5024091.pdf} \\
Tags            & \texttt{trs}, \texttt{bond}, \texttt{repo}, \texttt{par-spread},
                  \texttt{lgd}, \texttt{risky-annuity}, \texttt{practitioner} \\
Date read       & \today \\
\bottomrule
\end{tabular}

% =====================================================================
\section*{One-paragraph summary}
A practitioner write-up of bond TRS pricing in three components:
bond coupons, bond performance (change in price), and the funding
leg, all discounted with a \emph{risky discount factor}
$\widetilde P(t_0, t_i) = Q(t_0, t_i)\, P(t_0, t_i)$ that combines a
survival probability $Q$ and a risk-free DF $P$, plus a
loss-given-default (LGD) term to account for unwind on bond default.
The PV identity decomposes into four summations
(coupons + performance $-$ funding $-$ LGD), each carrying the same
risky DF. Future bond prices are projected at the repo rate,
$B_T = B_t(1 + r\,\tau_1 - r\,c\,\tau_2)$, with repo interest
deducted on intermediate coupons. The \emph{par spread} $p$ is
defined as the funding spread that makes PV zero; rearranging gives
$p = \mathrm{PV}(\mathrm{TRS}|s=0) / \mathrm{RiskyAnnuity}^{\mathrm{float}}$.
The appendix offers a compact alternative valuation: net the TRS as
``full bond cashflows to bond maturity $+$ funding leg to TRS
maturity $+$ a negative notional exchange equal to the current bond
price paid at TRS maturity'', useful for investment-grade bonds
where the LGD term is small.

% =====================================================================
\section{Setup and notation}

\begin{tabular}{@{}p{3cm}p{3.5cm}p{8.5cm}@{}}
\toprule
Paper                 & House                        & Meaning \\
\midrule
$\phi$                & $\phi \in \{+1, -1\}$        & $+1$ receiver of bond coupons, $-1$ payer \\
$N_B$                 & ---                          & Bond notional $=$ Units $\times$ Face Value \\
$N_C$                 & ---                          & Cash notional of the funding leg \\
$r$                   & ---                          & Bond coupon rate (fixed) \\
$\tau_i$              & $\alpha_i$                   & Year fraction of bond coupon period $i$ \\
$\tau_j$              & $\alpha_j$                   & Year fraction of funding-leg period $j$ \\
$F_j$                 & ---                          & Floating funding index fixing (e.g.\ SOFR) at $j$ \\
$s$                   & ---                          & Funding spread (added to $F$) \\
$p$                   & ---                          & Par / breakeven funding spread \\
$B(t)$                & $\Bond{t}{\cdot}$            & Dirty bond price at $t$ \\
$Q(t_0, t_i)$         & $Q(t_0, t_i)$                & Survival probability from $t_0$ to $t_i$ \\
$P(t_0, t_i)$         & $\Disc{t_0}{t_i}$            & Risk-free discount factor \\
$\widetilde P(t_0, t_i)$ & ---                       & Risky DF $= Q \cdot P$ \\
$RR$                  & ---                          & Recovery rate (paper convention: 40\% senior subordinated) \\
$1 - RR$              & ---                          & LGD per unit notional \\
$r^{\text{repo}}$     & $\Repo$                      & Repo rate (bond-specific) \\
$c$                   & ---                          & Coupon paid during projection window \\
$\tau_1, \tau_2$      & ---                          & Time from $t$ to $T$; time from coupon date to $T$ \\
\bottomrule
\end{tabular}

% =====================================================================
\section{Key equations}

\paragraph{TRS PV (top-level structure):}
\begin{equation}\label{eq:btrs-top}
  \mathrm{PV}_{\mathrm{TRS}} \;=\; \phi \Bigl( \mathrm{PV}(\text{Bond Coupons}) + \mathrm{PV}(\text{Performance}) - \mathrm{PV}(\text{Funding}) \Bigr),
\end{equation}
with $\phi = +1$ if receiving bond coupons (paying funding), $-1$
otherwise. LGD added separately as shown below.

\paragraph{PV of bond coupons.}
\begin{equation}\label{eq:btrs-coupons}
  \mathrm{PV}(\text{Bond Coupons}) \;=\; \sum_{i=1}^{n} N_B\, r\, \tau_i\, \underbrace{Q(t_0, t_i)\, P(t_0, t_i)}_{= \widetilde P(t_0, t_i)}.
\end{equation}
\textit{Says:} each fixed coupon $N_B r \tau_i$ is discounted by the
risky DF combining survival to $t_i$ and risk-free discounting from
$t_0$ to $t_i$.

\paragraph{PV of bond performance --- two settlement conventions.}
Performance paid on each coupon date:
\begin{equation}\label{eq:btrs-perf-periodic}
  \mathrm{PV}(\text{Performance}) \;=\; \sum_{i=1}^{n} N_B\, \bigl( B(t_{i-1}) - B(t_i) \bigr)\, Q(t_0, t_i)\, P(t_0, t_i).
\end{equation}
Performance paid at maturity $T$ only:
\begin{equation}\label{eq:btrs-perf-bullet}
  \mathrm{PV}(\text{Performance}) \;=\; N_B\, \bigl( B(t_0) - B(T) \bigr)\, Q(t_0, T)\, P(t_0, T).
\end{equation}
\textit{Says:} performance is the running change in bond price. The
periodic version reduces PFE and XVA costs (Trade Feature \#9) at the
expense of more frequent operational settlement.

\paragraph{PV of the funding leg.}
\begin{equation}\label{eq:btrs-funding}
  \mathrm{PV}(\text{Funding}) \;=\; \sum_{j=1}^{m} N_C\, (F_j + s)\, \tau_j\, Q(t_0, t_j)\, P(t_0, t_j).
\end{equation}
\textit{Says:} floating index $F$ plus spread $s$, accruing on cash
notional $N_C$. The same survival weighting applies because the TRS
unwinds on bond default --- the funding leg also stops.

\paragraph{Full TRS PV with LGD.}
\begin{equation}\label{eq:btrs-full}
  \begin{aligned}
  \mathrm{PV}_{\mathrm{TRS}} \;=\; \phi \biggl[\,
    & \sum_{i=1}^{n} N_B\, r\, \tau_i\, \widetilde P(t_0, t_i)
    + \sum_{i=1}^{n} N_B\, \bigl(B(t_{i-1}) - B(t_i)\bigr)\, \widetilde P(t_0, t_i) \\
    & - \sum_{j=1}^{m} N_C\, (F_j + s)\, \tau_j\, \widetilde P(t_0, t_j)
    \;-\; \sum_{i=1}^{n} N_B\, (1 - RR)\, \bigl(Q(t_0, t_{i-1}) - Q(t_0, t_i)\bigr)\, P(t_0, t_i)
  \,\biggr].
  \end{aligned}
\end{equation}
\textit{Says:} the fourth summation is the LGD term. Default in
$[t_{i-1}, t_i]$ has probability $Q(t_0, t_{i-1}) - Q(t_0, t_i)$;
multiplied by the (risk-free) discount and the LGD per unit notional.

\paragraph{Bond forward via repo} (Burgess uses simple compounding):
\begin{equation}\label{eq:btrs-fwd-noc}
  B_T \;=\; B_t\,(1 + r^{\text{repo}}\, \tau).
\end{equation}
With intermediate coupons $c$ deducted at the repo rate from coupon
date to $T$:
\begin{equation}\label{eq:btrs-fwd-coupon}
  B_T \;=\; B_t\,(1 + r^{\text{repo}}\, \tau_1) - c\,(1 + r^{\text{repo}}\, \tau_2)\, \big|_{\text{summed over coupons in } [t, T]}.
\end{equation}
\textit{Says:} the dirty bond is treated as ``on deposit'' earning
the repo rate. Intermediate coupons accrete to $T$ at the same rate.
This is the simple-compounding analogue of Pucci's
$\mathrm{Fwd}_t(t_e) = P^{mkt}_t e^{\Repo (t_e - t)} - \sum c_i e^{\Repo (t_e - t_i)}$.

\paragraph{Par (breakeven) funding spread.}
\begin{equation}\label{eq:btrs-par-spread}
  p \;=\; \frac{\mathrm{PV}(\mathrm{TRS}\,|\,s=0)}{\sum_j N_C\, \tau_j\, \widetilde P(t_0, t_j)},
  \qquad
  \text{Risky Annuity (float leg)} \;=\; \sum_j N_C\, \tau_j\, \widetilde P(t_0, t_j).
\end{equation}
\textit{Says:} solve PV $= 0$ for $s$; $p$ is the spread that makes
the two legs of equal value. The denominator is the risky annuity of
the funding leg.

\paragraph{Approximate decomposition} (Appendix):
\begin{itemize}[noitemsep]
  \item \textit{Full bond cashflows to bond maturity.}
  \item \textit{Funding leg to TRS maturity.}
  \item \textit{A negative notional exchange equal to the current
        bond price paid at TRS maturity} --- this is what represents
        the short performance leg in compact form.
  \item For investment-grade bonds, set $Q \equiv 1$ and drop LGD.
\end{itemize}

% =====================================================================
\section{Derivations \& commentary}

\begin{commentary}
The note is operational, not theoretical. The PV structure
\eqref{eq:btrs-full} treats every cashflow under the \emph{same}
risky DF $\widetilde P$, which is internally consistent given the
``TRS unwinds on default'' clause: if the bond defaults, both legs
stop simultaneously and an LGD payment closes out the position. The
single-curve setup hides several subtleties that the post-2008
Lou (2018) paper makes explicit (separate repo for the bond hedge
vs.\ collateral/OIS for derivative discounting).
\end{commentary}

\begin{commentary}
Two notational conflicts to watch when reading alongside other notes:
\textbf{(i)} Burgess uses $\phi$ for the long/short flag where Pucci
uses $a$ and others use $\theta$. \textbf{(ii)} $Q$ is the survival
probability here, not (as in some credit-derivatives texts) the
risk-neutral measure. The risky DF
$\widetilde P = Q \cdot P$ is then the credit-flavoured analogue of
Pucci's pre-default $\widehat{\Disc{t}{T}}$ under default-independence
(DI) and absolute-continuity (AC); they coincide under DI with
$Q(t,T) = e^{-\int_t^T \gamma_u\,\dd u}$.
\end{commentary}

\begin{commentary}
The repo-projection formula \eqref{eq:btrs-fwd-noc}--\eqref{eq:btrs-fwd-coupon}
is the simple-compounding version of what the anonymous Treasury
Lock Model and Pucci 2019 derive in continuous form. For short tenors
($\tau \lesssim 1$Y) the simple form is fine; for longer projections
the compound version (or the continuous $e^{r\tau}$ used in the
T-Lock paper) drifts apart and the choice matters.
\end{commentary}

\begin{commentary}
``Constant Units'' vs.\ ``Constant Notional'' is a contract feature,
not a pricing distinction. Under \emph{Constant Units} the bond
quantity is fixed and the cash notional varies (periodic notional
exchanges required, bond fixings needed). Under \emph{Constant
Notional} the cash side is fixed and the bond quantity drifts.
Pricing uses the same building blocks; only the bookkeeping of
$N_B$ vs.\ $N_C$ in successive periods changes.
\end{commentary}

\begin{commentary}
The appendix's compact valuation
(``bond + funding leg + negative notional payment at maturity'') is
useful because it lets one re-use existing bond-pricer and IRS-pricer
without building dedicated TRS code. For investment-grade names where
LGD is small it's accurate enough to be the production fallback. The
key insight is that the receiver of performance is essentially long
the bond \emph{minus} the bond price paid at TRS maturity, which is
exactly the payoff of buying the bond now and selling it at $T$ at
the forward price.
\end{commentary}

% =====================================================================
\section{Validation}

\begin{validation}
This is a \emph{specification} of mathematical objects the paper
defines and the numerical values they should take. It says nothing
about how those objects are implemented or invoked.

\textbf{Quantities the paper defines:}
\begin{itemize}[noitemsep]
  \item The risky discount factor $\widetilde P(t_0, t_i) =
        Q(t_0, t_i)\, P(t_0, t_i)$.
  \item Four PV components: bond coupons \eqref{eq:btrs-coupons},
        bond performance (periodic \eqref{eq:btrs-perf-periodic} or
        bullet \eqref{eq:btrs-perf-bullet}), funding leg
        \eqref{eq:btrs-funding}, and the LGD term in
        \eqref{eq:btrs-full}.
  \item Forward bond price under simple repo compounding
        \eqref{eq:btrs-fwd-noc}--\eqref{eq:btrs-fwd-coupon}.
  \item The par funding spread $p$ via \eqref{eq:btrs-par-spread} and
        the risky annuity of the float leg.
\end{itemize}

\textbf{Canonical numerical case} (pages 7--9):
\begin{itemize}[noitemsep]
  \item 50,000 bonds purchased at $104.5433\%$;
        cash equivalent USD $5{,}227{,}165$.
  \item Bond leg specified in Units.
  \item Funding: SOFR $+\,100$~bps.
  \item Performance paid at maturity only.
  \item Recovery rate $RR = 40\%$ (paper assumption for senior subordinated).
  \item Coupon worked example: $80{,}736 \times 100 \times 3.85\% \times 0.5 = 155{,}416.80$;
        $\text{Units} = 10{,}000{,}000 / (100 \times 123.86\%) = 80{,}736$.
\end{itemize}

\textbf{Expected numerical values}: no PV bottom-line printed in the
document --- only the per-period coupon arithmetic above. The
formulas are self-validating by structure.

\textbf{Equation $\leftrightarrow$ quantity map:}

\smallskip
\begin{tabular}{@{}p{2.8cm}p{3.8cm}p{6.2cm}@{}}
\toprule
Paper eq.                                 & Symbol                      & Quantity \\
\midrule
\eqref{eq:btrs-coupons}                   & $\mathrm{PV}(\text{Coupons})$       & risky-discounted bond coupon stream \\
\eqref{eq:btrs-perf-periodic}             & $\mathrm{PV}(\text{Performance})$   & periodic bond-price change PV \\
\eqref{eq:btrs-perf-bullet}               & $\mathrm{PV}(\text{Performance})$   & bullet (maturity-only) performance PV \\
\eqref{eq:btrs-funding}                   & $\mathrm{PV}(\text{Funding})$       & risky-discounted floating coupon PV \\
\eqref{eq:btrs-full}                      & $\mathrm{PV}_{\mathrm{TRS}}$        & full PV with LGD \\
\eqref{eq:btrs-fwd-noc}--\eqref{eq:btrs-fwd-coupon} & $B_T$              & forward bond price via simple repo \\
\eqref{eq:btrs-par-spread}                & $p$                          & breakeven funding spread \\
\bottomrule
\end{tabular}
\smallskip

\textbf{Invariants and identities to verify:}
\begin{enumerate}[noitemsep]
  \item \emph{Risky DF factorisation}: $\widetilde P(t_0, t_i) =
        Q(t_0, t_i)\, P(t_0, t_i)$ holds by definition; in
        particular $\widetilde P \to P$ as $Q \to 1$ (no default).
  \item \emph{Par spread reproduction}: pricing the TRS at spread
        $s = p$ from \eqref{eq:btrs-par-spread} returns
        $\mathrm{PV}_{\mathrm{TRS}} = 0$ within float tolerance.
  \item \emph{Periodic vs.\ bullet performance equivalence}: in the
        zero-default-zero-discounting limit ($Q \equiv 1$, $P \equiv 1$),
        $\sum_i (B(t_{i-1}) - B(t_i)) = B(t_0) - B(T)$ telescopes
        exactly; the two performance formulations differ only by
        survival/discount weights on intermediate periods.
  \item \emph{Investment-grade limit}: with $Q \equiv 1$, the LGD term
        vanishes (because the survival differences are zero) and
        \eqref{eq:btrs-full} reduces to a plain bond-vs-floater swap.
  \item \emph{Coupon arithmetic} (worked example): $80{,}736 \times
        100 \times 3.85\% \times 0.5 = 155{,}416.80$, and
        $10{,}000{,}000 / (100 \times 123.86\%) = 80{,}736$ (rounded).
  \item \emph{Default-period sum}: $\sum_{i=1}^{n}
        \bigl(Q(t_0, t_{i-1}) - Q(t_0, t_i)\bigr) = Q(t_0, t_0) -
        Q(t_0, t_n) = 1 - Q(t_0, t_n)$ (telescoping); used in the
        LGD term.
  \item \emph{Repo forward consistency}: \eqref{eq:btrs-fwd-noc} is
        the first-order Taylor approximation of $B_T = B_t e^{r\tau}$
        (Pucci's continuous form, eq.~50); the relative error at
        $r = 5\%$, $\tau = 1$Y is $\sim 0.12\%$. For TRS tenors
        beyond a few years, the compound form should be used.
\end{enumerate}
\end{validation}

% =====================================================================
\section{Glossary of symbols (paper's notation)}
\begin{tabular}{@{}p{3cm}p{12cm}@{}}
\toprule
$\phi$                & $+1$ receiver of bond / $-1$ payer of bond \\
$N_B$                 & Bond notional, $\text{Units} \times \text{Face Value}$ \\
$N_C$                 & Cash notional of the funding leg \\
$r$                   & Fixed bond coupon rate \\
$\tau_i$              & Year fraction of coupon period $i$ \\
$\tau_j$              & Year fraction of funding period $j$ \\
$F_j$                 & Floating funding index fixing at $j$ (e.g.\ SOFR) \\
$s$                   & Funding spread (added to $F$) \\
$p$                   & Par / breakeven funding spread \\
$B(t)$                & Dirty bond price at $t$ \\
$Q(t_0, t_i)$         & Survival probability from $t_0$ to $t_i$ \\
$P(t_0, t_i)$         & Risk-free discount factor from $t_0$ to $t_i$ \\
$\widetilde P(t_0, t_i)$ & Risky discount factor $= Q \cdot P$ \\
$RR$                  & Recovery rate (paper assumes 40\% senior subordinated) \\
$1 - RR$              & LGD per unit notional \\
$r^{\text{repo}}$     & Repo rate of the underlying bond \\
$c$                   & Coupon amount paid during a projection window \\
$\tau_1, \tau_2$      & Times: $t \to T$ and coupon date $\to T$ \\
$T$                   & TRS maturity date \\
\bottomrule
\end{tabular}

% =====================================================================
\begin{todobox}
\begin{itemize}[noitemsep]
  \item Build a worked end-to-end PV computation on the case-study
        inputs (50,000 units at $104.5433\%$, SOFR + 100~bps, $RR =
        40\%$, sensible curve). The paper sets up the inputs but
        does not print a bottom-line PV.
  \item Cross-check the simple-repo forward \eqref{eq:btrs-fwd-noc}
        against the continuous form $B_t e^{r\tau}$ for $\tau \in
        \{0.5, 1, 2, 5\}$Y at $r^{\text{repo}} \in \{2\%, 5\%\}$ ---
        quantify the gap.
  \item Implement the approximate decomposition in the appendix and
        compare with the full \eqref{eq:btrs-full} for an
        investment-grade name (low $1 - Q$); the gap should be the
        LGD term.
  \item Par-spread sensitivity to recovery rate at fixed curve and
        survival: should depend on $RR$ only through the LGD term.
  \item Reconcile the Burgess setup (single risky DF for both legs)
        with the Lou (2018) ``security leg + premium leg under
        different discounting'' picture. They should agree under
        Burgess's assumption that the TRS unwinds at default; they
        will disagree if the funding leg discounting is taken to be
        the dealer's funding rate.
\end{itemize}
\end{todobox}

\end{document}
